\chapter{电磁学：从规范不变性到 Maxwell 方程组}
\label{ch:electromagnetism}

\section{概述}

R3（局域 $U(1)$ 规范不变性）加上 R2（最小作用量原理）可以一条线推出全部经典电磁学。本章展示从规范原理到 Maxwell 四方程、Lorentz 力定律和电磁波方程的完整推导链。

核心洞察：\textbf{Maxwell 的四个方程不是独立的经验定律}——前两个（Gauss 磁定律、Faraday 定律）是场强张量 $F_{\mu\nu}$ 的 Bianchi 恒等式（几何必然），后两个（Gauss 电定律、Ampère-Maxwell 定律）是 $U(1)$ 规范作用量的 Euler-Lagrange 方程。，后两个（Gauss 电定律、Ampère-Maxwell 定律）是 $U(1)$ 规范作用量的 Euler-Lagrange 方程。

\section{Maxwell 方程组}

\subsection{Gauss 电定律}

\begin{lawbox}[Gauss 电定律]
\[
\nabla\cdot\mathbf{E} = \frac{\rho}{\varepsilon_0}
\]
\textbf{前提}：$\partial_\mu F^{\mu\nu} = J^\nu$ 的 $\nu=0$ 分量

\textbf{变量}：
\begin{itemize}
  \item $\mathbf{E}$ — 电场强度（V/m 或 N/C）。描述单位正电荷在该点受到的力
  \item $\rho$ — 电荷密度（C/m$^3$）。空间某点单位体积内的电荷量
  \item $\varepsilon_0 = 8.854\times 10^{-12}$ F/m — 真空介电常数。衡量真空对电场的"容纳"能力
  \item $\nabla\cdot\mathbf{E}$ — 电场的散度。正散度意味着电场线从该点"流出"（正电荷源），负散度意味着电场线"流入"（负电荷汇）
\end{itemize}
\end{lawbox}

\textbf{物理意义}：电荷是电场的"源"。电场线从正电荷发出，终止于负电荷。通过任意闭合曲面的电通量等于曲面内包含的总电荷除以 $\varepsilon_0$：$\oint_S \mathbf{E}\cdot d\mathbf{A} = Q_{\text{enc}}/\varepsilon_0$。

\subsection{Gauss 磁定律}

\begin{lawbox}[Gauss 磁定律]
\[
\nabla\cdot\mathbf{B} = 0
\]
\textbf{前提}：Bianchi 恒等式——$F_{\mu\nu} = \partial_\mu A_\nu - \partial_\nu A_\mu$ 的几何必然性

\textbf{变量}：
\begin{itemize}
  \item $\mathbf{B}$ — 磁感应强度（Tesla, T = N/(A·m)）。描述运动电荷在磁场中受到的力
  \item $\nabla\cdot\mathbf{B} = 0$ — 磁场散度恒为零。这意味着磁力线没有起点和终点
\end{itemize}
\end{lawbox}

\textbf{物理意义}：不存在磁单极子（magnetic monopole）。磁力线总是闭合的回线——进入一个区域多少，就流出多少。任何切割磁铁都会产生新的南北极对，而不是孤立的北极或南极。这是 $F_{\mu\nu} = \partial_\mu A_\nu - \partial_\nu A_\mu$ 这种"反对称导数"结构的几何必然——从数学上禁止了 $\nabla\cdot\mathbf{B} \neq 0$。

\subsection{Faraday 电磁感应定律}

\begin{lawbox}[Faraday 感应定律]
\[
\nabla\times\mathbf{E} = -\frac{\partial\mathbf{B}}{\partial t}
\]
\textbf{前提}：Bianchi 恒等式——与 Gauss 磁定律同源

\textbf{变量}：
\begin{itemize}
  \item $\nabla\times\mathbf{E}$ — 电场的旋度。衡量电场在空间中的"旋转"程度
  \item $\partial\mathbf{B}/\partial t$ — 磁场随时间的变化率。负号反映了 Lenz 定律：感应电流的方向总是\textbf{抵抗}引起它的磁通量变化
\end{itemize}
\end{lawbox}

\textbf{物理意义}：变化的磁场产生（有旋的）电场。这是发电机和变压器的工作原理。将一个磁铁推入线圈，变化的磁通量产生感应电动势——线圈中就有了电流。

\subsection{Ampère-Maxwell 定律}

\begin{lawbox}[Ampère-Maxwell 定律]
\[
\nabla\times\mathbf{B} = \mu_0\mathbf{J} + \mu_0\varepsilon_0\frac{\partial\mathbf{E}}{\partial t}
\]
\textbf{前提}：$\partial_\mu F^{\mu\nu} = J^\nu$（Euler-Lagrange 方程，需要 R2 做变分）

\textbf{变量}：
\begin{itemize}
  \item $\mathbf{J}$ — 电流密度（A/m$^2$）。单位面积流过的电流
  \item $\mu_0 = 4\pi\times 10^{-7}$ N/A$^2$ — 真空磁导率。衡量真空对磁场的"传导"能力
  \item $\mu_0\varepsilon_0\partial\mathbf{E}/\partial t$ — \textbf{位移电流}（displacement current）。Maxwell 为理论自洽性引入的关键修正项。它的物理来源是：变化的电场也产生磁场，就像真实的电流一样
  \item 常数关系：$c = 1/\sqrt{\mu_0\varepsilon_0} = 2.998\times 10^8$ m/s —— 光速从纯电磁常数中算出，是 Maxwell 最伟大的成就
\end{itemize}
\end{lawbox}

\begin{derivationbox}[为什么需要位移电流项？]
  Maxwell 之前的 Ampère 定律只有 $\nabla\times\mathbf{B} = \mu_0\mathbf{J}$。但取散度得：
  \[
  \nabla\cdot(\nabla\times\mathbf{B}) = \mu_0\nabla\cdot\mathbf{J} = 0
  \]
  这意味着 $\nabla\cdot\mathbf{J} = 0$（稳恒电流）。但在非稳恒情况下（如电容器充放电），电荷守恒要求：
  \[
  \nabla\cdot\mathbf{J} + \frac{\partial\rho}{\partial t} = 0
  \]
  Maxwell 意识到需要在 Ampère 定律中加入 $\partial\mathbf{E}/\partial t$ 项。利用 Gauss 定律 $\nabla\cdot\mathbf{E} = \rho/\varepsilon_0$，有 $\partial\rho/\partial t = \varepsilon_0\nabla\cdot(\partial\mathbf{E}/\partial t)$。因此修正后的方程：
  \[
  \nabla\times\mathbf{B} = \mu_0\mathbf{J} + \mu_0\varepsilon_0\frac{\partial\mathbf{E}}{\partial t}
  \]
  取散度得 $0 = \mu_0(\nabla\cdot\mathbf{J} + \varepsilon_0\nabla\cdot(\partial\mathbf{E}/\partial t)) = \mu_0(\nabla\cdot\mathbf{J} + \partial\rho/\partial t)$——与电荷守恒自洽！
\end{derivationbox}

\section{Lorentz 力定律}

\begin{lawbox}[Lorentz 力]
\[
\mathbf{F} = q(\mathbf{E} + \mathbf{v}\times\mathbf{B})
\]
\textbf{前提}：$U(1)$ 规范场与带电粒子的最小耦合（minimal coupling）

\textbf{变量}：
\begin{itemize}
  \item $q$ — 粒子电荷（C）。正电荷受力与 $\mathbf{E}$ 同向，负电荷反向
  \item $\mathbf{v}$ — 粒子速度（m/s）
  \item $\mathbf{v}\times\mathbf{B}$ — 磁力部分。磁力总是垂直于速度——\textbf{磁场不做功}（$\mathbf{F}_B\cdot\mathbf{v}=0$），只能改变运动方向
  \item $\mathbf{E} + \mathbf{v}\times\mathbf{B}$ — 电磁力的完整表达式。在粒子的静止系中磁力消失
\end{itemize}
\end{lawbox}

\begin{derivationbox}[推导：最小耦合原理]
  在 $U(1)$ 规范理论中，带电粒子与电磁场的相互作用通过"最小耦合"实现：将经典动量替换为规范协变动量：
  \[
  \mathbf{p} \;\to\; \mathbf{p} - q\mathbf{A}
  \]
  其中 $\mathbf{A}$ 是磁矢势（$\mathbf{B} = \nabla\times\mathbf{A}$），标量势 $\phi$ 满足 $\mathbf{E} = -\nabla\phi - \partial\mathbf{A}/\partial t$。

  拉格朗日量变为 $L = \frac{1}{2}m\dot{\mathbf{r}}^2 - q\phi + q\dot{\mathbf{r}}\cdot\mathbf{A}$。Euler-Lagrange 方程给出：
  \[
  \mathbf{F} = q(-\nabla\phi - \partial_t\mathbf{A}) + q\mathbf{v}\times(\nabla\times\mathbf{A})
  = q\mathbf{E} + q\mathbf{v}\times\mathbf{B}
  \]
\end{derivationbox}

\section{电磁波方程}

\begin{lawbox}[电磁波方程]
\[
\nabla^2\mathbf{E} = \frac{1}{c^2}\frac{\partial^2\mathbf{E}}{\partial t^2},\qquad
\nabla^2\mathbf{B} = \frac{1}{c^2}\frac{\partial^2\mathbf{B}}{\partial t^2}
\]
\[
c = \frac{1}{\sqrt{\varepsilon_0\mu_0}} \approx 2.998\times 10^8\ \text{m/s}
\]
\textbf{前提}：Maxwell 方程组的真空解（$\rho=0$, $\mathbf{J}=0$）

\textbf{波动方程的标准形式}为 $\nabla^2\psi = (1/v^2)\partial^2\psi/\partial t^2$，其中 $v$ 是波速。从纯电磁常数 $\varepsilon_0,\mu_0$ 算出的 $c$ 与实测光速精确一致——\textbf{光就是电磁波}。
\end{lawbox}

\begin{derivationbox}[推导：从 Maxwell 到波动方程]
  \textbf{第 1 步：取 Faraday 定律的旋度。}
  \[
  \nabla\times(\nabla\times\mathbf{E}) = -\nabla\times\frac{\partial\mathbf{B}}{\partial t} = -\frac{\partial}{\partial t}(\nabla\times\mathbf{B})
  \]

  \textbf{第 2 步：用矢量恒等式展开左边，用 Ampère-Maxwell 代入右边。}
  \[
  \underbrace{\nabla(\nabla\cdot\mathbf{E})}_{\text{=0 (真空中}\rho=0\text{)}} - \nabla^2\mathbf{E}
  = -\frac{\partial}{\partial t}\!\left(\mu_0\mathbf{J} + \mu_0\varepsilon_0\frac{\partial\mathbf{E}}{\partial t}\right)
  \]

  \textbf{第 3 步：整理得标准波动方程。}
  \[
  \nabla^2\mathbf{E} = \mu_0\varepsilon_0\frac{\partial^2\mathbf{E}}{\partial t^2} = \frac{1}{c^2}\frac{\partial^2\mathbf{E}}{\partial t^2}
  \]
  对 $\mathbf{B}$ 同法可得完全相同的波动方程，波速同为 $c = 1/\sqrt{\mu_0\varepsilon_0}$。

  \textbf{物理意义}：Maxwell 在 1865 年推导出这个方程时，已知的 $\varepsilon_0$ 和 $\mu_0$ 值代入后得到的波速 $c \approx 3.1\times 10^8$ m/s——与 Fizeau 1849 年测得的光速 $3.15\times 10^8$ m/s 在实验误差内一致。他由此断言"光本身（包括辐射热）是一种电磁扰动"。这是物理学史上最伟大的统一之一。
\end{derivationbox}

\section{推导关系总结}

\begin{center}
\begin{tikzpicture}[
  node distance=0.8cm,
  kernel/.style={rectangle,draw,fill=principleblue!15,rounded corners,text width=2cm,align=center,font=\footnotesize\bfseries},
  law/.style={rectangle,draw,fill=lawgreen!15,rounded corners,text width=2.6cm,align=center,font=\footnotesize},
  arrow/.style={->,>=Stealth,thick,deriveorange}
]
  \node[kernel] (r3) {R3\\U(1) 规范};
  \node[kernel,right=0.3cm of r3] (r2) {R2\\$\delta S = 0$};
  \node[law,below=0.5cm of r3] (ge) {Gauss 电\\$\nabla\cdot\mathbf{E}=\rho/\varepsilon_0$};
  \node[law,below=0.3cm of ge] (gm) {Gauss 磁\\$\nabla\cdot\mathbf{B}=0$};
  \node[law,below=0.3cm of gm] (fa) {Faraday\\$\nabla\times\mathbf{E}=-\partial_t\mathbf{B}$};
  \node[law,below=0.3cm of fa] (am) {Ampère-Maxwell\\$\nabla\times\mathbf{B}=\mu_0\mathbf{J}+\mu_0\varepsilon_0\partial_t\mathbf{E}$};
  \node[law,below=0.3cm of am] (lf) {Lorentz 力\\$\mathbf{F}=q(\mathbf{E}+\mathbf{v}\times\mathbf{B})$};
  \node[law,below=0.3cm of lf] (wave) {电磁波 $\nabla^2\mathbf{E}=c^{-2}\partial_t^2\mathbf{E}$\\$c=1/\sqrt{\varepsilon_0\mu_0}$};
  \draw[arrow] (r3) -- (ge); \draw[arrow] (r3) -- (gm);
  \draw[arrow] (gm) -- (fa); \draw[arrow] (r2) -- (am);
  \draw[arrow] (ge) -- (am); \draw[arrow] (fa) -- (am);
  \draw[arrow] (r3) -- (lf); \draw[arrow] (am) -- (wave); \draw[arrow] (fa) -- (wave);
\end{tikzpicture}
\end{center}

参考文献：[GRIFF] §7--9, [JACK] §5--7, [FEYN II] Ch.\ 18--20, Maxwell (1865)
